\exercisesection{}

\begin{enumerate}
\def\labelenumi{\arabic{enumi})}
\tightlist
\item
  设 \(E = \ell_{\mathbf{C}}^{2}(\mathbf{N})\)，令 \(u \in \mathcal{L}(\mathrm{E})\) 为满足 \(u(x_{0}, \ldots, x_{n}, \ldots) = (0, x_{0}, \ldots, x_{n}, \ldots)\) 的自同态，其中该关系对 E 中所有 \((x_{n})_{n \in \mathbf{N}}\) 都成立。令 \((e_{n})\) 为 E 的标准正交基。
\end{enumerate}

\begin{enumerate}
\def\labelenumi{\alph{enumi})}
\item
  对每个 \(v \in \mathcal{L}(\mathrm{E})\) 满足 \(\|v\| < 1/2\)，向量 \(e_{0}\) 不在 \(u + v\) 的像的闭包中。
\item
  \(\mathcal{L}(\mathrm{E})\) 中以 u 为中心、半径为 1/2 的开球不包含 E 的任何可对角化自同态。
\end{enumerate}

\begin{enumerate}
\def\labelenumi{\arabic{enumi})}
\setcounter{enumi}{1}
\item
  设 E 是复希尔伯特空间。直接证明 E 的自同态 \(u\) 是弗雷德霍姆自同态，当且仅当 \(\pi(u)\) 在 \(\mathcal{C}alk(E)\) 中可逆（阿特金森定理）。（使用紧算子的谱性质，参见 IV 第 146 页第 1 节。）
\item
  设 E 是复希尔伯特空间，\(\mathrm{B}=(e_{i})_{i\in\mathrm{I}}\) 是 E 的标准正交基。对每个 \(i\in I\)，记 \(p_{i}\) 为像为 \(Ce_{i}\) 的正交投影。
\end{enumerate}

\begin{enumerate}
\def\labelenumi{\alph{enumi})}
\item
  设 K 是 E 的相对紧子集。对每个 \(\varepsilon > 0\)，存在 I 的有限子集 J，使得对每个 \(x \in \mathbf{K}\)，都有 \(\|x - \sum_{i \in \mathbf{J}} p_j(x)\| \leqslant \varepsilon\)。
\item
  设 \(u \in \mathcal{D}_{\mathrm{B}}(\mathrm{E})\)，并记其特征值族为 \(\lambda = (\lambda_{i})_{i \in \mathrm{I}}\)。族 \((\lambda_{i} p_{i})\) 在赋予相对紧收敛拓扑的 \(\mathcal{L}(\mathrm{E})\) 中可求和，且其和等于 u（参见 IV 第 147 页命题 1）。
\end{enumerate}

\begin{enumerate}
\def\labelenumi{\arabic{enumi})}
\setcounter{enumi}{3}
\item
  设 E 是无限维复希尔伯特空间。存在 E 的不是有限迹的紧自同态。
\item
  \begin{enumerate}
  \def\labelenumii{\alph{enumii})}
  \tightlist
  \item
  设 \(E = L^{2}([0,1], \mu)\)，其中 \(\mu\) 是勒贝格测度。代数 \(L^{\infty}([0,1], \mu)\) 在如下映射下的像 A 是 \(\mathcal{L}(E)\) 的极大交换子代数：该映射把 f 对应到乘以 f 的自同态；A 不具有某个 E 的标准正交基 B 所对应的 \(\mathcal{D}_{\mathrm{B}}(\mathrm{E})\) 形式。（验证代数 A 不含极小投影。）
  \end{enumerate}
\end{enumerate}

\begin{enumerate}
\def\labelenumi{\alph{enumi})}
\setcounter{enumi}{1}
\tightlist
\item
  设 E 是无限维希尔伯特空间。存在 \(\mathcal{L}(\mathrm{E})\) 的极大交换子代数，它不是对于 E 的某个标准正交基 B 而具有 \(\mathcal{D}_{\mathrm{B}}(\mathrm{E})\) 的形式。
\end{enumerate}

\begin{enumerate}
\def\labelenumi{\arabic{enumi})}
\setcounter{enumi}{5}
\tightlist
\item
  设 E 是复希尔伯特空间。
\end{enumerate}

\begin{enumerate}
\def\labelenumi{\alph{enumi})}
\item
  设 u 是 E 的紧正规自同态。E 的每个在 u 下不变的闭子空间也在 \(u^{*}\) 下不变。
\item
  若 E 无限维，则存在 E 的正规自同态 u 和闭子空间 \(F \subset E\)，使得：
\end{enumerate}

\begin{enumerate}
\def\labelenumi{(\roman{enumi})}
\item
  有 \(u(\mathrm{F}) \subset \mathrm{F}\)；
\item
  空间 F 不是 \(u^{*}\)-不变的；
\item
  空间 F° 不是在 u 下不变的；
\item
  由 \(u\) 通过限制到子空间在 F 上诱导的自同态不是正规的。
\end{enumerate}

\begin{enumerate}
\def\labelenumi{\arabic{enumi})}
\setcounter{enumi}{6}
\tightlist
\item
  本习题的目标是给出 IV 第 149 页定理 1 的一个更直接的证明。设 E 是无限维复希尔伯特空间，u 是 E 的紧正规自同态。
\end{enumerate}

\begin{enumerate}
\def\labelenumi{\alph{enumi})}
\item
  证明，对每个非零复数 \(\lambda\)，特征空间 \(\operatorname{Ker}(u - \lambda 1_{\mathrm{E}})\) 是有限维的。
\item
  假设 \(u\) 是埃尔米特的。证明 E 中存在 \(x_0 \neq 0\)，使得
\end{enumerate}

\[
\frac {\| u (x _ {0}) \| ^ {2}}{\| x _ {0} \| ^ {2}} = \sup _ {x \in \mathrm{E} - \{0 \}} \frac {\| u (x) \| ^ {2}}{\| x \| ^ {2}}.
\]

\begin{enumerate}
\def\labelenumi{\alph{enumi})}
\setcounter{enumi}{2}
\item
  证明 \(x_{0}\) 是 \(u\) 关于某个特征值 \(\lambda\) 的特征向量，且 \(|\lambda| = \|u\|\)。
\item
  证明存在一个可数集 I、E 中的标准正交族 \((e_{i})_{i\in\mathrm{I}}\) 以及实数族 \((\lambda_{i})_{i\in\mathrm{I}}\)，使得
\end{enumerate}

\[
u (x) = \sum_ {i \in I} \lambda_ {i} \langle e _ {i} | x \rangle e _ {i}
\]

对所有 \(x \in E\) 都成立。

\begin{enumerate}
\def\labelenumi{\alph{enumi})}
\setcounter{enumi}{4}
\tightlist
\item
  将结论推广到 u 紧且正规的情形。（写成 \(u = u_{1} + iu_{2}\)，其中 \(u_{1}\) 和 \(u_{2}\) 是紧的埃尔米特自同态且彼此可交换。）
\end{enumerate}

\begin{enumerate}
\def\labelenumi{\arabic{enumi})}
\setcounter{enumi}{7}
\tightlist
\item
  我们赋予群 R/Z 以其归一化的哈尔测度 \(\mu\)。
\end{enumerate}

\begin{enumerate}
\def\labelenumi{\alph{enumi})}
\tightlist
\item
  设 \(h \in \mathcal{L}^{2}(\mathbf{R}/\mathbf{Z})\)。对 \(f \in \mathcal{L}^{2}(\mathbf{R}/\mathbf{Z})\)，按下式在 R/Z 上定义函数 \(\mathrm{C}_{h}(f)\)：
\end{enumerate}

\[
\mathrm{C} _ {h} (f) (x) = \int_ {\mathbf {R} / \mathbf {Z}} f (t) h (x - t) d \mu (t).
\]

映射 \(f \mapsto \mathrm{C}_h(f)\) 是 \(\mathcal{L}^2(\mathbf{R}/\mathbf{Z})\) 的连续自同态，并通过取商定义 \(\mathrm{L}^2(\mathbf{R}/\mathbf{Z})\) 的紧正规自同态。

\begin{enumerate}
\def\labelenumi{\alph{enumi})}
\setcounter{enumi}{1}
\item
  确定 \(C_{h}\) 的谱以及该谱中各元素的谱重数。（将 h 展开为傅里叶级数。）
\item
  存在 \(h \in \mathcal{C}(\mathbf{R}/\mathbf{Z})\)，使得 \(C_{h}\) 不是有限迹的。
\end{enumerate}

\begin{enumerate}
\def\labelenumi{\arabic{enumi})}
\setcounter{enumi}{8}
\tightlist
\item
  记 E 为复希尔伯特空间 \(L^{2}([0,1])\)，其元素是在 \([0,1]\) 上关于勒贝格测度 \(\mu\) 平方可积的函数。
\end{enumerate}

令 \(u_{0} \in \mathcal{L}(E)\) 为 E 上满足下述条件的沃尔泰拉自同态：

\[
u _ {0} (f) (x) = \int_ {[ 0, x ]} f (t) d \mu (t)
\]

对所有 \(f \in \mathrm{L}^2([0,1])\) 以及几乎所有 \(x \in [0,1]\) 都成立。

\begin{enumerate}
\def\labelenumi{\alph{enumi})}
\item
  自同态 \(u_{0}\) 是紧的且 \(\operatorname{Sp}(u_{0})=\{0\}\)。
\item
  自同态 \(u_0^* u_0\) 具有有限迹；计算其迹并由此求出级数
\end{enumerate}

\[
\sum_ {n = 1} ^ {+ \infty} \frac {1}{n ^ {2}}.
\]

\begin{enumerate}
\def\labelenumi{\alph{enumi})}
\setcounter{enumi}{2}
\tightlist
\item
  令 \(v_{0}\) 为 E 上由下式定义的秩 1 自同态
\end{enumerate}

\[
v _ {0} (f) = \left(\int_ {0} ^ {1} f (t) d \mu (t)\right) \cdot 1,
\]

其中 1 表示常值函数 1 的等价类。令 \(u = u_0 - v_0\)。证明 \(u\) 是希尔伯特–施密特算子，并计算满足下式的核 N：

\[
u (f) (x) = \int_ {0} ^ {1} \mathrm{N} (x, y) f (y) d \mu (y)
\]

对 \(f \in \mathrm{E}\) 和几乎所有 \(x \in [0,1]\) 都成立。

\begin{enumerate}
\def\labelenumi{\alph{enumi})}
\setcounter{enumi}{3}
\item
  确定 u 的特征值及其谱重数。（计算 \(u(e_{n})\)，其中 \(e_{n}(t) = e^{2i\pi nt}\) 对所有 \(n \in Z\) 都成立。）
\item
  计算 \(u^{*}u\) 的迹，并再次求出通项为 \(1/n^{2}\)（\(n \geqslant 1\)）的级数的值。
\item
  将这些论证推广，证明若 \(k \geqslant 1\) 是严格正整数，则
\end{enumerate}

\[
\sum_ {n = 1} ^ {+ \infty} \frac {1}{n ^ {2 k}} \in \pi^ {2 k} {\bf Q} ^ {*}.
\]

\begin{enumerate}
\def\labelenumi{\arabic{enumi})}
\setcounter{enumi}{9}
\tightlist
\item
  设 E 是复希尔伯特空间。令 A 为 \(\mathcal{L}(\mathrm{E})\) 的星子代数。假设 A 在 E 中的表示是不可约的，且 A 含有非零紧自同态。
\end{enumerate}

\begin{enumerate}
\def\labelenumi{\alph{enumi})}
\item
  星子代数 A 含有非零有限秩投影。（证明 A 含有非零紧埃尔米特自同态，并应用函数演算。）
\item
  星子代数 A 含有秩 1 投影。（若 \(p \in \mathrm{A}\) 是最小秩的非零投影，证明代数 \(p\mathrm{A}p\) 包含于 \(\mathbf{C}p\)。）
\item
  从而 A 含有 \(\mathcal{L}^{\mathrm{c}}(\mathrm{E})\)。
\item
  反之，\(\mathcal{L}(\mathrm{E})\) 的每个含有 \(\mathcal{L}^{\mathrm{c}}(\mathrm{E})\) 的星子代数都是不可约的。
\end{enumerate}

\begin{enumerate}
\def\labelenumi{\arabic{enumi})}
\setcounter{enumi}{10}
\tightlist
\item
  设 E 是复希尔伯特空间。令 \(A \subset \mathcal{L}(\mathrm{E})\) 为紧自同态的集合，使其满足 \(u \circ v = v \circ u\)：A 中每一对元素 \((u, v)\) 都满足此关系，且 \(u^{*} \in A\) 对所有 \(u \in A\) 成立。
\end{enumerate}

存在一个可数标准正交族 \((e_i)_{i\in \mathrm{I}}\) 以及一个映射族 \((\Lambda_i)_{i\in \mathrm{I}}\)，其元素都是 \(\mathrm{A}\to \mathbf{C}\) 的映射，使得

\[
u (x) = \sum_ {i \in \mathrm{I}} \Lambda_ {i} (u) \langle e _ {i} | x \rangle e _ {i}
\]

对 A 中的每个元素 u 和 E 中的每个 x 都有。若 A 是 \(\mathcal{L}(\mathrm{E})\) 的含单位元子代数，则 \(\Lambda_{i}\) 对每个 \(i \in I\) 都是 A 的一个特征。

\begin{enumerate}
\def\labelenumi{\arabic{enumi})}
\setcounter{enumi}{11}
\tightlist
\item
  若 E 是复希尔伯特空间且 u 是 E 的紧正规自同态，证明 u 有循环向量（参见 IV 第 191 页定义 3）当且仅当满足以下条件：
\end{enumerate}

\begin{enumerate}
\def\labelenumi{(\roman{enumi})}
\item
  u 的核为 0；
\item
  u 的每个非零特征值的谱重数都等于 1。
\end{enumerate}

\begin{enumerate}
\def\labelenumi{\arabic{enumi})}
\setcounter{enumi}{12}
\tightlist
\item
  设 E 是可数型复希尔伯特空间，\((e_{n})_{n\in\mathbf{N}}\) 是 E 的标准正交基，u 是 E 的自同态。证明 u 紧，当且仅当
\end{enumerate}

\[
\lim _ {n \to + \infty} \sup _ {x \in \{e _ {0}, \dots , e _ {n - 1} \} ^ {\circ} - \{0 \}} \frac {\| u (x) \| ^ {2}}{\| x \| ^ {2}} = 0.
\]

\begin{enumerate}
\def\labelenumi{\arabic{enumi})}
\setcounter{enumi}{13}
\item
  令 \(\mathrm{X} = [0,1]\) 赋以勒贝格测度，并令 \(\mathrm{N}\colon \mathrm{X}\times \mathrm{X}\to \mathbf{R}_{+}\) 满足 \(\mathrm{N}(x,y) = |x - y|\)。自同态 \(u_{\mathrm{N}}\)（作用于 \(\mathrm{L}^2 (\mathrm{X})\)）不是正的。
\item
  设 X 是局部紧拓扑空间，\(\mu\) 是 X 上的正有界测度且支撑为 X，并令 \(N \in \mathcal{C}_{b}(X \times X)\)。
\end{enumerate}

\begin{enumerate}
\def\labelenumi{\alph{enumi})}
\tightlist
\item
  存在一个自同态 \(u\)，其作用于 \(\mathcal{C}_b(\mathrm{X})\)，使得
\end{enumerate}

\[
u (f) (y) = \int_ {\mathrm{X}} \mathrm{N} (x, y) f (x) d \mu (x)
\]

对每个函数 \(f \in \mathcal{C}_b(\mathrm{X})\) 和每个 \(y \in \mathrm{X}\) 都成立。下文假设 \(u\) 是紧的；例如当空间 \(\mathrm{X}\) 紧时就是如此。

\begin{enumerate}
\def\labelenumi{\alph{enumi})}
\setcounter{enumi}{1}
\item
  \(u_{\mathrm{N}}\) 的像包含于 \(\mathcal{C}_b(\mathrm{X})\) 中，且 \(u_{\mathrm{N}}\) 定义了从 \(\mathrm{L}_{\mathbf{C}}^{2}(\mathrm{X})\) 到 \(\mathcal{C}_b(\mathrm{X})\) 的连续线性映射。
\item
  设 \(f \in \mathrm{L}_{\mathbf{C}}^{2}(\mathrm{X})\) 且 \(\lambda \in \mathbf{C}^*\)，满足 \(u_{\mathrm{N}}(f) = \lambda f\)。则 \(f \in \mathcal{C}_b(\mathrm{X})\)。
\end{enumerate}

\(d)\) 存在一个可数族 \((f_i)_{i\in I}\)，其元素属于 \(\mathcal{C}_b(\mathrm{X})\)，以及一个非零实数族 \((\lambda_i)_{i\in I}\)，使得对所有 \(f\in \mathrm{L}_{\mathbf{C}}^2 (\mathrm{X})\)，有

\[
u (f) = \sum_ {i \in I} \lambda_ {i} \langle f _ {i} | f \rangle f _ {i}
\]

在 \(\mathcal{C}_b(\mathrm{X})\) 中。

\begin{enumerate}
\def\labelenumi{\alph{enumi})}
\setcounter{enumi}{4}
\tightlist
\item
  假设 \(u_{\mathrm{N}}\) 是 \(\mathrm{L}_{\mathbf{C}}^{2}(\mathrm{X})\) 的正自同态。对每个 \(i \in \mathrm{I}\)，记 \(h_i \in \mathcal{C}_b(\mathrm{X} \times \mathrm{X})\) 为由 \(h_i(x,y) = \overline{f_i(x)} f_i(y)\) 定义的函数。则
\end{enumerate}

\[
N = \sum_ {i \in I} \lambda_ {i} h _ {i}
\]

在赋予紧收敛拓扑的 \(\mathcal{C}(\mathrm{X}\times\mathrm{X})\) 中。（使用迪尼定理（TG，第 X 卷，第 34 页，定理 1）证明该级数收敛，然后证明等式。）

\begin{enumerate}
\def\labelenumi{\arabic{enumi})}
\setcounter{enumi}{15}
\tightlist
\item
  令 \(\mathrm{X} = [0,1]\) 赋以勒贝格测度。
\end{enumerate}

\begin{enumerate}
\def\labelenumi{\alph{enumi})}
\item
  对序列 \((f_n)_{n\in \mathbf{N}}\)（其元素是 X 上的连续函数，且在 \(\mathrm{L}_{\mathbf{C}}^2 (\mathrm{X})\) 中标准正交），存在核 \(k\in \mathcal{C}(\mathrm{X}\times \mathrm{X})\)，使得
\item
  自同态 \(u_{k}\)（作用于 \(\mathrm{L}_{\mathbf{C}}^{2}(\mathrm{X})\)）是正的；
\item
  对每个 \(n \in \mathbf{N}\)，函数 \(f_n\) 是 \(u_k\) 的特征函数，其特征值非零；
\item
  族 \((f_n)\) 是自同态 \(u_k\) 的像的一个基。
\item
  存在核 \(k \in \mathcal{C}(\mathrm{X} \times \mathrm{X})\) 以及特征函数标准正交族 \((f_n)\)，这些是自同态 \(u_k\)（作用于 \(\mathrm{L}_{\mathbf{C}}^2(\mathrm{X})\)）的有界特征函数，使得序列 \((\|f_n\|_\infty)\) 无界。
\end{enumerate}

\begin{enumerate}
\def\labelenumi{\arabic{enumi})}
\setcounter{enumi}{16}
\item
  设 E、F 是希尔伯特空间。双线性映射 \((u,v)\mapsto\) \(\mathrm{Tr}(uv)\) 从巴拿赫空间 \(\mathcal{L}_1(\mathrm{E};\mathrm{F})'\) 到 \(\mathcal{L}(\mathrm{F};\mathrm{E})\) 定义了一个等距同构。
\item
  设 E 是复希尔伯特空间，令 \(I_{E}\) 为 IV 第 151 页第 3 号中定义的集合。
\end{enumerate}

令 \(n \in I_{E}\)。令 F 为复希尔伯特空间。映射 \(u \mapsto \lambda_{n}(|u|)\) 从 \(\mathcal{L}^{c}(E;F)\) 到 \(R_{+}\) 是连续的。

\begin{enumerate}
\def\labelenumi{\arabic{enumi})}
\setcounter{enumi}{18}
\item
  保持 IV 第 161 页引理 5 的记号。证明若 \(f\) 的像包含于 \(\mathbf{R}\)，且 \(f\) 严格凸、\(\varrho_n > 0\)，则不等式 (6) 中取等号，当且仅当 \(a_i = b_i\) 对 \(0 \leqslant i \leqslant n\) 成立。
\item
  设 E 是预希尔伯特空间，F 是巴拿赫空间，且 \(u: E \to F\) 是线性映射。
\end{enumerate}

\begin{enumerate}
\def\labelenumi{\alph{enumi})}
\item
  证明 u 是紧线性映射，当且仅当 \(\|u(e_{n})\|\to0\) 对 E 中的每个标准正交序列 \((e_{n})\) 都成立。（为证明充分性，使用 TG，第 IX 卷，第 20 页命题 14。）
\item
  假设 E 是可数型希尔伯特空间且 F=E。存在 E 的一个标准正交基 \((e_{n})_{n\in\mathbf{N}}\)，使 \(\|u(e_{n})\|\to0\)，其充要条件是 u 不是弗雷德霍姆自同态。
\item
  假设 E 是希尔伯特空间且 F=E。若 \(\langle u(e_{n}) | e_{n} \rangle \to 0\) 对 E 的每个标准正交序列 \((e_{n})\) 都成立，则 u 是紧线性映射。
\end{enumerate}

\begin{enumerate}
\def\labelenumi{\arabic{enumi})}
\setcounter{enumi}{20}
\tightlist
\item
  设 E 是可数型希尔伯特空间。
\end{enumerate}

\begin{enumerate}
\def\labelenumi{\alph{enumi})}
\item
  设 A 是 E 中一个沿有限集补集滤子弱收敛于 0 的无限单位向量集。证明存在一个取值于 A 的序列 \((x_{n})_{n\geqslant1}\) 和 E 中的标准正交序列 \((e_{n})_{n\geqslant1}\)，使得 \(\|x_{n}-e_{n}\|\leqslant1/n\) 对所有 \(n\geqslant1\) 都成立。
\item
  设 u 是从 E 到 E 的线性映射，不必连续。假设对 E 中每个标准正交序列 \((e_{n})\)，序列 \((u(e_{n}))\) 都弱收敛于 0。证明 u 连续且紧。
\item
  设 u 是 E 的自同态。存在 E 的一个标准正交基 \((e_{n})_{n\in\mathbf{N}}\)，使序列 \((u(e_{n}))_{n\in\mathbf{N}}\) 弱收敛于 0，当且仅当不存在 E 的自同态 v，使得 \(v\circ u-1_{E}\) 是紧的。
\end{enumerate}

\begin{enumerate}
\def\labelenumi{\arabic{enumi})}
\setcounter{enumi}{21}
\tightlist
\item
  设 E 是复希尔伯特空间。令 u 和 v 为 E 的正自同态。
\end{enumerate}

\begin{enumerate}
\def\labelenumi{\alph{enumi})}
\tightlist
\item
  对每个满足 \(0 < p < 1\) 的实数 p 以及每个 \(x \in R_{+}\)，有
\end{enumerate}

\[
x ^ {p} = \frac {\sin (p \pi)}{\pi} \int_ {0} ^ {+ \infty} t ^ {p} \left(\frac {1}{t} - \frac {1}{t + x}\right) d t.
\]

\begin{enumerate}
\def\labelenumi{\alph{enumi})}
\setcounter{enumi}{1}
\item
  假设 \(u \leqslant v\)。对每个实数 \(p\)，满足 \(0 < p < 1\)，有 \(u^p \leqslant v^p\)。（使用 I 第 119 页引理 14 的 b) 和上一问。）\\
\item
  令 \(n \in \mathbf{N}\)。置 \(a = v^{n/2} \circ u^n \circ v^{n/2}\)，且 \(b = (v^{1/2} \circ u \circ v^{1/2})^n\)。若 \(a \leqslant 1_{\mathrm{E}}\)，则 \(b \leqslant 1_{\mathrm{E}}\)。
\item
  假设 u 和 v 具有有限迹。对每个整数 \(n \in N\)，有
\end{enumerate}

\[
\operatorname{Tr} \left(\left(v ^ {1 / 2} \circ u \circ v ^ {1 / 2}\right) ^ {n}\right) \leqslant \operatorname{Tr} \left(v ^ {n / 2} \circ u ^ {n} \circ v ^ {n / 2}\right).
\]

（令 \(a\) 和 \(b\) 如问题 \(c\) 所定义；利用 IV 第 161 页引理 5，把问题归结为证明 \(\widehat{\Lambda}^{n}b\) 的最大特征值小于 \(\widehat{\Lambda}^{n}a\) 的最大特征值；以 \(\widehat{\Lambda}^{n}\mathrm{E}\) 替代 E，并以 \(u\)、\(v\) 分别替换为 \(\widehat{\Lambda}^{n}u\)、\(\widehat{\Lambda}^{n}v\)，然后应用上一问。）
% semantic-fragments: legacy-input-seam
\relax{}
